% Options for packages loaded elsewhere
% Options for packages loaded elsewhere
\PassOptionsToPackage{unicode}{hyperref}
\PassOptionsToPackage{hyphens}{url}
%
\documentclass[
  ignorenonframetext,
  aspectratio=169,
  russian,
]{beamer}
\newif\ifbibliography
\usepackage{pgfpages}
\setbeamertemplate{caption}[numbered]
\setbeamertemplate{caption label separator}{: }
\setbeamercolor{caption name}{fg=normal text.fg}
\beamertemplatenavigationsymbolshorizontal
% remove section numbering
\setbeamertemplate{part page}{
  \centering
  \begin{beamercolorbox}[sep=16pt,center]{part title}
    \usebeamerfont{part title}\insertpart\par
  \end{beamercolorbox}
}
\setbeamertemplate{section page}{
  \centering
  \begin{beamercolorbox}[sep=12pt,center]{section title}
    \usebeamerfont{section title}\insertsection\par
  \end{beamercolorbox}
}
\setbeamertemplate{subsection page}{
  \centering
  \begin{beamercolorbox}[sep=8pt,center]{subsection title}
    \usebeamerfont{subsection title}\insertsubsection\par
  \end{beamercolorbox}
}
% Prevent slide breaks in the middle of a paragraph
\widowpenalties 1 10000
\raggedbottom
\AtBeginPart{
  \frame{\partpage}
}
\AtBeginSection{
  \ifbibliography
  \else
    \frame{\sectionpage}
  \fi
}
\AtBeginSubsection{
  \frame{\subsectionpage}
}
\usepackage{iftex}
\ifPDFTeX
  \usepackage[T1]{fontenc}
  \usepackage[utf8]{inputenc}
  \usepackage{textcomp} % provide euro and other symbols
\else % if luatex or xetex
  \usepackage{unicode-math} % this also loads fontspec
  \defaultfontfeatures{Scale=MatchLowercase}
  \defaultfontfeatures[\rmfamily]{Ligatures=TeX,Scale=1}
\fi
\usepackage{lmodern}

\ifPDFTeX\else
  % xetex/luatex font selection
\fi
% Use upquote if available, for straight quotes in verbatim environments
\IfFileExists{upquote.sty}{\usepackage{upquote}}{}
\IfFileExists{microtype.sty}{% use microtype if available
  \usepackage[]{microtype}
  \UseMicrotypeSet[protrusion]{basicmath} % disable protrusion for tt fonts
}{}


\usepackage{longtable,booktabs,array}
\usepackage{calc} % for calculating minipage widths
\usepackage{caption}
% Make caption package work with longtable
\makeatletter
\def\fnum@table{\tablename~\thetable}
\makeatother
\usepackage{graphicx}
\makeatletter
\newsavebox\pandoc@box
\newcommand*\pandocbounded[1]{% scales image to fit in text height/width
  \sbox\pandoc@box{#1}%
  \Gscale@div\@tempa{\textheight}{\dimexpr\ht\pandoc@box+\dp\pandoc@box\relax}%
  \Gscale@div\@tempb{\linewidth}{\wd\pandoc@box}%
  \ifdim\@tempb\p@<\@tempa\p@\let\@tempa\@tempb\fi% select the smaller of both
  \ifdim\@tempa\p@<\p@\scalebox{\@tempa}{\usebox\pandoc@box}%
  \else\usebox{\pandoc@box}%
  \fi%
}
% Set default figure placement to htbp
\def\fps@figure{htbp}
\makeatother



\ifLuaTeX
\usepackage[bidi=basic,provide=*]{babel}
\else
\usepackage[bidi=default,provide=*]{babel}
\fi
% get rid of language-specific shorthands (see #6817):
\let\LanguageShortHands\languageshorthands
\def\languageshorthands#1{}


\setlength{\emergencystretch}{3em} % prevent overfull lines

\providecommand{\tightlist}{%
  \setlength{\itemsep}{0pt}\setlength{\parskip}{0pt}}



 

\usepackage[]{csquotes}

\IfFileExists{plex-otf.sty}{
  %% Full TeXlive
  % \usepackage[%
  %   % math,
  %   RM={Scale=0.94},SS={Scale=0.94},SScon={Scale=0.94},TT={Scale=MatchLowercase,FakeStretch=0.9},DefaultFeatures={Ligatures=Common}
  % ]{plex-otf}
}{
  %% TinyTeX
  \usepackage{libertine}
}

%%% Load theme
% https://deic.uab.cat/~iblanes/beamer_gallery/
\IfFileExists{beamerthemegotham.sty}{
  %% Full TeXlive
  \usetheme{gotham}
  \gothamset{
    numbering=totalpagenumber,
    parttocframe default=off,
    sectiontocframe default=off,
    subsectiontocframe default=off,
  }
}{
  %% TinyTeX
  \usetheme{Madrid}
}
\makeatletter
\@ifpackageloaded{caption}{}{\usepackage{caption}}
\AtBeginDocument{%
\ifdefined\contentsname
  \renewcommand*\contentsname{Содержание}
\else
  \newcommand\contentsname{Содержание}
\fi
\ifdefined\listfigurename
  \renewcommand*\listfigurename{Список иллюстраций}
\else
  \newcommand\listfigurename{Список иллюстраций}
\fi
\ifdefined\listtablename
  \renewcommand*\listtablename{Список таблиц}
\else
  \newcommand\listtablename{Список таблиц}
\fi
\ifdefined\figurename
  \renewcommand*\figurename{Рисунок}
\else
  \newcommand\figurename{Рисунок}
\fi
\ifdefined\tablename
  \renewcommand*\tablename{Таблица}
\else
  \newcommand\tablename{Таблица}
\fi
}
\@ifpackageloaded{float}{}{\usepackage{float}}
\floatstyle{ruled}
\@ifundefined{c@chapter}{\newfloat{codelisting}{h}{lop}}{\newfloat{codelisting}{h}{lop}[chapter]}
\floatname{codelisting}{Список}
\newcommand*\listoflistings{\listof{codelisting}{Листинги}}
\makeatother
\makeatletter
\makeatother
\makeatletter
\@ifpackageloaded{caption}{}{\usepackage{caption}}
\@ifpackageloaded{subcaption}{}{\usepackage{subcaption}}
\makeatother

\usepackage{bookmark}
\IfFileExists{xurl.sty}{\usepackage{xurl}}{} % add URL line breaks if available
\urlstyle{same}
\hypersetup{
  pdftitle={Моделирование неравновесной агрегации},
  pdfauthor={Казазаев Даниил Михайлович; Симонова Виктория Игоревна; Малюга Валерия Васильевна; Наталья какая-то; Лисенков Егор Романович},
  pdflang={ru-RU},
  hidelinks,
  pdfcreator={LaTeX via pandoc}}


\title{Моделирование неравновесной агрегации}
\subtitle{DLA и фрактальные структуры}
\author{Казазаев Даниил Михайлович \and Симонова Виктория
Игоревна \and Малюга Валерия Васильевна \and Наталья
какая-то \and Лисенков Егор Романович}
\date{2026-03-21}

\begin{document}
\frame{\titlepage}


\begin{frame}{Введение: Неравновесная агрегация}
\phantomsection\label{ux432ux432ux435ux434ux435ux43dux438ux435-ux43dux435ux440ux430ux432ux43dux43eux432ux435ux441ux43dux430ux44f-ux430ux433ux440ux435ux433ux430ux446ux438ux44f}
\begin{itemize}[<+->]
\tightlist
\item
  \textbf{Неравновесная агрегация} - ключевой процесс во многих
  природных и технических явлениях:

  \begin{itemize}[<+->]
  \tightlist
  \item
    образование сажи,
  \item
    рост дендритов при электроосаждении,
  \item
    «вязкие пальцы» при вытеснении нефти,
  \item
    электрический пробой.
  \end{itemize}
\item
  Общая черта: \textbf{необратимое прилипание} частиц к растущему
  кластеру (\(E_{\text{bond}} \gg k_B T\)).
\item
  Результат - \textbf{разветвлённые, часто фрактальные структуры}.
\end{itemize}
\end{frame}

\begin{frame}{Модель DLA (Witten-Sander)}
\phantomsection\label{ux43cux43eux434ux435ux43bux44c-dla-witten-sander}
\begin{block}{Основные принципы}
\phantomsection\label{ux43eux441ux43dux43eux432ux43dux44bux435-ux43fux440ux438ux43dux446ux438ux43fux44b}
\begin{itemize}[<+->]
\tightlist
\item
  Пространство - \textbf{квадратная решётка}.
\item
  В центре - \textbf{затравочная частица} (seed).
\item
  Новые частицы \textbf{диффундируют} (случайные блуждания) извне.
\item
  При контакте с кластером - \textbf{необратимое прилипание}.
\end{itemize}
\end{block}

\begin{block}{Алгоритм}
\phantomsection\label{ux430ux43bux433ux43eux440ux438ux442ux43c}
\begin{enumerate}[<+->]
\tightlist
\item
  \textbf{Запуск частицы} с круга радиусом
  \(R_{\text{spawn}} = R_{\max} + \Delta R\)\\
  \(\alpha = 2\pi \cdot \text{random}\)
\item
  \textbf{Случайное блуждание} (4 направления, равные вероятности).
\item
  \textbf{Прилипание} - если соседний узел занят.
\item
  Если частица уходит за \(R_{\text{kill}} \approx 3R_{\max}\) -
  уничтожить.
\item
  Повторять до нужного числа частиц.
\end{enumerate}
\end{block}
\end{frame}

\begin{frame}{Морфология роста: Эффект экранирования}
\phantomsection\label{ux43cux43eux440ux444ux43eux43bux43eux433ux438ux44f-ux440ux43eux441ux442ux430-ux44dux444ux444ux435ux43aux442-ux44dux43aux440ux430ux43dux438ux440ux43eux432ux430ux43dux438ux44f}
\textbf{Почему структура ветвится, а не образует сплошной диск?} 1.
\textbf{Геометрическое преимущество:} Выступающие части кластера
(«ветви») имеют более высокую вероятность столкновения с блуждающей
частицей. 2. \textbf{Экранирование впадин:} Внутренние области (фьорды)
защищены внешними ветвями; вероятность проникновения частицы вглубь
экспоненциально мала. 3. \textbf{Математическая связь:} Процесс
описывается \textbf{уравнением Лапласа} \(\nabla^2 P = 0\), где
вероятность прилипания пропорциональна градиенту поля на границе.
\end{frame}

\begin{frame}{Фрактальная размерность}
\phantomsection\label{ux444ux440ux430ux43aux442ux430ux43bux44cux43dux430ux44f-ux440ux430ux437ux43cux435ux440ux43dux43eux441ux442ux44c}
\begin{block}{Что такое фрактал?}
\phantomsection\label{ux447ux442ux43e-ux442ux430ux43aux43eux435-ux444ux440ux430ux43aux442ux430ux43b}
\begin{itemize}[<+->]
\tightlist
\item
  Объект с \textbf{дробной размерностью}.
\item
  Плотность убывает с ростом размера: \[
  \rho(R) \sim R^{D-d}
  \]
\item
  Для DLA на плоскости \(D \approx 1.71 \pm 0.02\).
\end{itemize}
\end{block}
\end{frame}

\begin{frame}{Фрактальная размерность}
\phantomsection\label{ux444ux440ux430ux43aux442ux430ux43bux44cux43dux430ux44f-ux440ux430ux437ux43cux435ux440ux43dux43eux441ux442ux44c-1}
\begin{block}{Методы определения \(D\)}
\phantomsection\label{ux43cux435ux442ux43eux434ux44b-ux43eux43fux440ux435ux434ux435ux43bux435ux43dux438ux44f-d}
\begin{longtable}[]{@{}ll@{}}
\caption{Основные методы расчёта фрактальной
размерности}\label{tbl-methods}\tabularnewline
\toprule\noalign{}
Метод & Принцип \\
\midrule\noalign{}
\endfirsthead
\toprule\noalign{}
Метод & Принцип \\
\midrule\noalign{}
\endhead
\textbf{Сфер} & \(m(R) \sim R^D\) → наклон \(\ln m\) от \(\ln R\) \\
\textbf{Box counting} & \(N(L) \sim L^{-D}\) \\
\textbf{Радиус гирации} & \(N \sim R_g^D\) \\
\bottomrule\noalign{}
\end{longtable}
\end{block}
\end{frame}

\begin{frame}{Методы анализа и верификации}
\phantomsection\label{ux43cux435ux442ux43eux434ux44b-ux430ux43dux430ux43bux438ux437ux430-ux438-ux432ux435ux440ux438ux444ux438ux43aux430ux446ux438ux438}
\begin{itemize}[<+->]
\item
  \textbf{Метод радиуса гирации (\(R_g\)):}\\
  \[
  R_g = \sqrt{\frac{1}{N} \sum_{i=1}^{N} (\vec{r}_i - \vec{r}_{cm})^2}
  \]\\
  Построение \(\ln N\) от \(\ln R_g\) даёт наклон \(D\).
\item
  \textbf{Метод Box-counting:}\\
  Покрытие кластера сеткой ячеек \(\epsilon\), подсчёт непустых ячеек
  \(N(\epsilon)\).\\
  \[
  D = \lim_{\epsilon \to 0} \frac{\ln N(\epsilon)}{\ln(1/\epsilon)}
  \]
\end{itemize}
\end{frame}

\begin{frame}{Примеры математических фракталов}
\phantomsection\label{ux43fux440ux438ux43cux435ux440ux44b-ux43cux430ux442ux435ux43cux430ux442ux438ux447ux435ux441ux43aux438ux445-ux444ux440ux430ux43aux442ux430ux43bux43eux432}
\begin{longtable}[]{@{}ll@{}}
\caption{Классические фракталы и их
размерности}\label{tbl-classical-fractals}\tabularnewline
\toprule\noalign{}
Фрактал & Размерность \\
\midrule\noalign{}
\endfirsthead
\toprule\noalign{}
Фрактал & Размерность \\
\midrule\noalign{}
\endhead
Множество Кантора & \(\ln 2 / \ln 3 \approx 0.631\) \\
Кривая Коха & \(\ln 4 / \ln 3 \approx 1.262\) \\
Треугольник Серпинского & \(\ln 3 / \ln 2 \approx 1.585\) \\
Траектория броуновской частицы & \(2\) (без самопересечений) \\
\bottomrule\noalign{}
\end{longtable}
\end{frame}

\begin{frame}{Модификации модели}
\phantomsection\label{ux43cux43eux434ux438ux444ux438ux43aux430ux446ux438ux438-ux43cux43eux434ux435ux43bux438}
\begin{block}{Химически ограниченная агрегация (RLA)}
\phantomsection\label{ux445ux438ux43cux438ux447ux435ux441ux43aux438-ux43eux433ux440ux430ux43dux438ux447ux435ux43dux43dux430ux44f-ux430ux433ux440ux435ux433ux430ux446ux438ux44f-rla}
\begin{itemize}[<+->]
\tightlist
\item
  Вероятность прилипания \(p < 1\).
\item
  Может зависеть от числа соседей: \[
  p_1 = 0.01,\; p_2 = 0.1,\; p_3 = 0.9
  \]
\item
  Размерность увеличивается (\(D \to 2\)).
\end{itemize}
\end{block}

\begin{block}{Баллистическая агрегация}
\phantomsection\label{ux431ux430ux43bux43bux438ux441ux442ux438ux447ux435ux441ux43aux430ux44f-ux430ux433ux440ux435ux433ux430ux446ux438ux44f}
\begin{itemize}[<+->]
\tightlist
\item
  Частицы движутся \textbf{прямолинейно}.
\item
  Кластер более плотный, но граница - фрактал (\(D \approx 1.95\)).
\end{itemize}
\end{block}
\end{frame}

\begin{frame}{Модификации модели}
\phantomsection\label{ux43cux43eux434ux438ux444ux438ux43aux430ux446ux438ux438-ux43cux43eux434ux435ux43bux438-1}
\begin{block}{Кластер-кластерная агрегация (DLCA)}
\phantomsection\label{ux43aux43bux430ux441ux442ux435ux440-ux43aux43bux430ux441ux442ux435ux440ux43dux430ux44f-ux430ux433ux440ux435ux433ux430ux446ux438ux44f-dlca}
\begin{itemize}[<+->]
\tightlist
\item
  Одновременный рост множества кластеров.
\item
  Кластеры могут слипаться.
\item
  Размерность \textbf{меньше}, чем в DLA.
\end{itemize}
\end{block}
\end{frame}

\begin{frame}{Связь с электрическим пробоем}
\phantomsection\label{ux441ux432ux44fux437ux44c-ux441-ux44dux43bux435ux43aux442ux440ux438ux447ux435ux441ux43aux438ux43c-ux43fux440ux43eux431ux43eux435ux43c}
\begin{itemize}[<+->]
\tightlist
\item
  Модель DLA - частный случай \textbf{модели диэлектрического пробоя}
  (DBM).
\item
  Вероятность ветвления пропорциональна локальной напряжённости поля.
\item
  При показателе \(\eta = 1\) - эквивалент DLA.
\end{itemize}
\end{frame}

\begin{frame}{Заключение}
\phantomsection\label{ux437ux430ux43aux43bux44eux447ux435ux43dux438ux435}
\begin{itemize}[<+->]
\tightlist
\item
  DLA - простая, но мощная модель \textbf{неравновесного роста}.
\item
  Позволяет исследовать \textbf{фрактальные свойства} кластеров.
\item
  Модификации расширяют применимость к реальным системам.
\item
  Вычислительная реализация даёт наглядное представление о
  \textbf{скейлинговых закономерностях}.
\end{itemize}
\end{frame}




\end{document}
